\ifx\danslelivre\undefined
\documentclass[10pt]{../fichiers-configuration-latex/smfart}
\usepackage{ifthen}
\newboolean{english}
\newboolean{francais}
\setboolean{english}{false}
\setboolean{francais}{true}
\input{../fichiers-configuration-latex/paquets}
\selectlanguage{francais}
\input{../fichiers-configuration-latex/lemme-proposition-theoreme-ascii}
\input{../fichiers-configuration-latex/numerotation}
\input{../fichiers-configuration-latex/macros}

\externaldocument{../00-introduction/introduction}
\externaldocument{../02-topologies/topologies}
\externaldocument{../03-approximation/approximation}
\externaldocument{../04-Cohen-Gabber/Cohen-Gabber}
\externaldocument{../05-algebrisation-partielle/algebrisation}
\externaldocument{../06-log-lissite/log}
\externaldocument{../07-uniformisation-locale/uniformisation-faible}
%\externaldocument{../08-modification-theorem/modification}
\externaldocument{../09-uniformisation-premiere-a-l/uniformisation-locale-premiere-a-l}
\externaldocument{../11-produits-orientes/produits-orientes}
\externaldocument{../12-descente-cohomologique-orientee/descente-orientee}
\externaldocument{../13-theoreme-de-finitude/finitude}
\externaldocument{../14-fonctions-dimension/fonctions-dimension}
\externaldocument{../15-Lefschetz-affine/Lefschetz-affine}
\externaldocument{../16-classes-de-Chern/gysin}
\externaldocument{../17-complexes-dualisants/dualite}
\externaldocument{../19-exemples-contre-exemples/contre-exemples}
\externaldocument{../20-rigidite/rigidite}
\externaldocument{../21-finitude-non-abelienne/finitude-non-abelienne}

\begin{document}

\begin{center}
I. Anneaux excellents\\
Michel Raynaud rédigé par Yves Laszlo
\end{center}

\tableofcontents

\else
\setboolean{english}{false}
\setboolean{francais}{true}
\selectlanguage{francais}
\chapter{Anneaux excellents}
\addtocontents{toc}{par Michel Raynaud, rédigé par Yves Laszlo}
\begin{center}
Michel Raynaud, rédigé par Yves Laszlo
\end{center}
\fi

\section{Introduction} Ce texte est une version  un peu modifiée d'un exposé donné par Michel
Raynaud au printemps 2006 dans le cadre d'un séminaire sur les travaux de Gabber sur la cohomologie
des schémas excellents. Le but est de familiariser le lecteur avec la notion d'excellence et de lui
donner un fil d'Ariane pour se repérer dans \textbf{ÉGA IV} où l'on trouve les principales
propriétés des anneaux excellents. Son ambition n'est certainement pas de donner une exposition
complète de la théorie, mais une idée de la stratégie qui ramène pour l'essentiel les preuves à des
énoncés, souvent difficiles, dans le cas complet. Dans un second temps, on montre que toutes les
propriétés définissant les anneaux excellents peuvent être mises en défaut, même en petite
dimension. Notamment, il existe des anneaux de valuation discrète non excellents ainsi que des
anneaux noethériens intègres de dimension $1$ dont le lieu régulier n'est pas ouvert. Ce dernier
exemple est un sous-produit d'une construction proposée par Gabber (\refext{XIX}{19-rem-cexgabber}).
Elle montre que le théorème de constructibilité des images directes
(\refext{XIII}{13-constructibilité}) n'est plus vrai si on omet la condition de quasi-excellence.



\section{Définitions} Soit $A$ un anneau noethérien et $X=\Spec(A)$ son spectre. On va s'intéresser
à des conditions sur $X$ de deux sortes.

\textbf{$\bullet$ Conditions globales :}

\subsection{Condition 1 : conditions d'ouverture} \emph{Tout schéma intègre $Y$ fini sur $X$
contient un ouvert dense \begin{enumerate} \item[1.a)]  régulier.  \item[1.b)] normal.
\end{enumerate}}


\begin{remarque}\label{11->ouvert}La condition 1.a)  entraîne d'après le critère d'ouverture de
Nagata que le lieu régulier de tout schéma fini sur $X$ est ouvert (\egalong{IV}{2}{6.12.4}). De même,
la condition 1.b) entraîne que le lieu normal de tout schéma fini sur $X$ est ouvert
(\egalong{IV}{2}{6.13.7})\footnote{Et en fait, 1.a) (resp. 1.b)) entraîne que le lieu régulier (resp.
normal) de tout schéma intègre de type fini sur $X$ est ouvert}. Ces critères d'ouverture  assurent
en outre que pour tester 1.a) ou 1.b) on peut se limiter à des schémas intègres $Y$  qui sont de
plus finis \emph{et} radiciels sur $X$.  \end{remarque}

\textbf{$\bullet$  Conditions locales.}

Elles sont de deux types.

\subsection{Condition 2 : Conditions sur les fibres formelles} \emph{Pour tout point \emph{fermé}
$x$ de $X$, le morphisme de complétion \footnotemark $\Spec(\chap{\OO_x})\ra\Spec(\OO_x)$ est
\begin{enumerate} \item[2.a)] régulier.  \item[2.b)] normal.  \item[2.c)] réduit.  \end{enumerate} }
\footnotetext{Ses fibres sont appelées les fibres formelles (de $X$ ou $A$) en $x$.}

Un anneau vérifiant 2.a) est dit « G-ring » en anglais, ce en l'honneur de Grothendieck qui a
dégagé l'importance de la notion et étudié ses propriétés.

\begin{remarque}\label{1rem-fibre} Rappelons qu'un morphisme de schémas noethériens est dit régulier
(resp. normal, réduit) respectivement s'il est plat et si les fibres géométriques en tout point de
la base sont régulières (resp. normales, réduites). On dit que les fibres formelles de $X$ en $x$
sont géométriquement régulières, géométriquement normales ou géométriquement réduites si le
morphisme de complétion $\Spec(\chap{\OO_x})\ra\Spec(\OO_x)$ est régulier, normal ou réduit. Bien
entendu, il suffit de tester la régularité, normalité, ou réduction des fibres après changement de
base radiciel fini (\egalong{IV}{2}{6.7.7}).  Notons que la fibre fermée de
$\Spec(\chap{\OO_x})\ra\Spec(\OO_x)$ est le spectre du corps résiduel $k(x)$ : elle est toujours
géométriquement régulière. La fibre formelle en $y\in\Spec(\OO_x)$ s'identifie à la fibre formelle
\emph{générique} du sous-schéma fermé $\overline{\{y\}}$ (muni de sa structure réduite), adhérence
de $y$ dans $\Spec(\OO_{x})$ ; ceci explique qu'on s'intéresse dans la littérature aux fibres
formelles génériques des anneaux intègres. Dans le cas où $A$ est local mais pas un corps, elles
peuvent avoir des dimensions arbitraires entre $0$ et $\dim(A)-1$ et contenir des points fermés de
hauteurs différentes, même dans le cas excellent (\ref{1def-excellent}) régulier (\cite{Rot91}).
Dans le cas où $A$ est un localisé d'une algèbre intègre de type fini sur un corps, la dimension de
la fibre formelle générique est bien $\dim(A)-1$ (\cite{Mat88}).  \end{remarque}




\subsection{Condition 3 : condition de formelle caténarité} \emph{Pour tout point fermé $y$ d'un
sous-schéma fermé irréductible $Y$ de $X$, le complété\footnotemark $\chap{\OO_{Y,y}}$ est
équidimensionnel.  }\footnotetext{Bien entendu, même si $\OO_{Y,y}$ est intègre, son complété n'est
en général pas intègre : penser à une courbe nodale.}

On dit alors que $X$ est formellement caténaire. Par exemple, si $X$ est de dimension $1$,  $X$ est
formellement caténaire.

\begin{exemple}\label{1reg-form-cat}\label{1complet-formel-cat} Tout anneau local noethérien
\emph{complet} est formellement caténaire.\end{exemple}


\bigskip

Rappelons que $X$ est dit \emph{caténaire} si toutes les chaînes saturées de fermés irréductibles
de $X$ ayant mêmes extrémités ont même longueur, \emph{universellement
caténaire}\footnote{\label{1form-dim}Cette dernière notion est utile en théorie de la dimension :
si $A$ est intègre universellement caténaire contenue dans $B$ intègre de type fini sur $A$, on a
pour tout $\p\in\Spec(B)$ au dessus de $\q\in\Spec(A)$ la formule $$\dim
B_\p+\text{degtr}_{k(\q)}k(\p)=\dim A_\q+\text{degtr}_AB.$$ Mais, en pratique, on teste plutôt la
formelle caténarité qui, comme on le voit juste après, entraîne l'universelle caténarité, et même
lui est équivalente (voir~(\ref{1Ratliff}) plus bas) !} si tout schéma affine de type fini sur $X$
est caténaire. La caténarité est une notion locale. La terminologie de \emph{formelle caténarité}
est alors justifiée par la proposition élémentaire suivante (\egalong{IV}{2}{7.1.4}), proposition qui
résulte de la fidèle platitude du morphisme de complétion

\begin{lemme}\label{1form-versus-cat} \label{1note-cat-non uni cat}Soit $A$ local noethérien de
complété équidimensionnel. Alors \begin{itemize} \item[i)] $A$ est équidimensionnel et caténaire.
\item[ii)]Pour tout idéal $I$ de $A$, le quotient $A/I$ est équidimensionnel si et seulement si son
complété l'est ; en particulier, $A/I$ est formellement caténaire.  \item[iii)] En particulier, un
schéma affine $X$ noethérien formellement caténaire est caténaire et même universellement caténaire.
\end{itemize} \end{lemme}

Notons que \textit{iii)} découle immédiatement de \textit{i)} puisque $X$ est caténaire si et
seulement si ses composantes irréductibles le sont. On verra plus bas dans la section~\ref{1inv-exc}
que la propriété de formelle caténarité est notamment stable par extension finie d'où l'universelle
caténarité annoncée (cf. la preuve de la proposition~\ref{1equi-excellent} et, pour une réciproque,
voir (\ref{1Ratliff})).

\begin{exemple}\label{1CM-form-cat}Soit alors $B\ra A$ un morphisme local surjectif d'anneaux
noethériens et supposons $B$ de Cohen-Macaulay (par exemple régulier).  Comme $\chap{B}$ est de
Cohen-Macaulay, il est équidimensionnel de sorte que $A$ est formellement caténaire d'après
(\ref{1form-versus-cat}).  \end{exemple}

Regardons ce qui se passe dans le cas complet. Rappelons pour mémoire le théorème de structure de
Cohen des anneaux locaux complets noethériens (\egalong{IV}{1}{0.19.8.8}):

\begin{theoreme}[Cohen]\label{1Co} Soit $A$ un anneau local noethérien complet de corps résiduel
$k$.\begin{itemize} \item [(i)] $A$ est isomorphe à un quotient d'un anneau de séries formelles sur
un anneau de Cohen\footnote{Rappelons (\egalong{IV}{2}{19.8.5}) que les anneaux de Cohen $C$ sont les
corps de caractéristique nulle et les anneaux de valuation discrète complets d'inégale
caractéristique $p$ non ramifiés.  Lorsque le corps résiduel $\kappa$ de $C$ est parfait, $C$ n'est
autre que l'anneau des vecteurs de Witt de $\kappa$.}. Si $A$ contient  un corps, il est isomorphe à
un quotient d'un anneau de séries formelles sur $k$.  \item [(ii)] Si $A$ est de plus intègre, il
existe un sous-anneau $B$ isomorphe à un anneau de séries formelles sur un anneau de Cohen ou un
corps\footnote{Voir (\ref{1Coh-Gab}) pour une amélioration.} de sorte que l'inclusion $B\ra A$ soit
locale,  finie et  induise un isomorphisme  des corps résiduels.  \end{itemize} \end{theoreme}

Tout anneau local noethérien complet est donc quotient d'un anneau régulier.

\begin{definition}\label{1def-excellent} Soit $X$ un schéma (resp. $X=\Spec(A)$ un schéma affine)
noethérien. On dit que $X$ (resp. $A$) est \begin{itemize} \item \emph{excellent} si $X$ vérifie
$1.a)+2.a)+3)$.  \item \emph{quasi-excellent} si $X$ vérifie $1.a)+2.a)$.  \item
\emph{universellement japonais}\footnote{Ou \emph{Nagata} en anglais, voire
\emph{pseudo-géométrique} (chez Nagata notamment).} si $X$ vérifie 1.b)+2.c).  \end{itemize}
\end{definition}

\subsection{}\label{1desing} L'existence d'une classe de schémas stable par extension finie pour
laquelle le théorème de désingularisation est vérifié impose de se limiter aux schémas
quasi-excellents. Précisément, si tous les schémas intègres et finis $Y$ sur $X$ admettent une
désingularisation (au sens de l'existence de $Y'\ra Y$ propre et birationnel avec $Y'$ régulier),
alors $X$ est quasi-excellent \egalong{IV}{2}{7.9.5})\footnote{Si de plus $X$ peut localement se
plonger dans un schéma régulier, alors $X$ vérifie 3) et est donc excellent.}.  Inversement, le
théorème de désingularisation d'Hironaka se généralise à tout schéma réduit quasi-excellent de
caractéristique nulle (\cite[3.4.3]{desingularization@Temkin})\footnote{Ce résultat a été longtemps
considéré comme « bien connu des experts » alors que sa preuve, tout à fait non triviale, date
de 2008.}

On regroupe plus bas (\ref{1pathologie}) des exemples de « méchants anneaux ». Commençons par un
regard plus positif.

\section{Exemples immédiats.}

\begin{proposition} Un corps, un anneau de Dedekind de corps des fractions \emph{de caractéristique
nulle} est excellent.\end{proposition}

\begin{proof} Vérifions qu'un corps est excellent. En effet, une algèbre finie et intègre sur un
corps est un corps : les propriétés 1.a), 2.a) et 3) sont donc vérifiées ce qui prouve que tout
corps est excellent.

\medskip

Soit $A$ un anneau de Dedekind de corps des fractions $K$ \emph{de caractéristique nulle} est
excellent.

\begin{itemize} \item Vérifions 1.a). Soit donc $B$ intègre finie sur $A$. Soit $B$ est un corps,
auquel cas 1.a) est vérifié, soit $A$ se plonge dans $B$. Comme $K$ est de caractéristique nulle,
$B$ est génériquement étale sur $A$, prouvant que le lieu régulier de $B$ contient un ouvert non
vide (le lieu étale par exemple).  \item Pour 2.a), considérons $x$ fermé dans $\Spec(A)$. La fibre
formelle non fermée en $x$ est le complété $\chap{ K_x}$ de $K$ pour la valuation définie par $x$.
Comme $K$ est de caractéristique nulle, le corps $\chap{ K_x}$ est séparable sur $K$ d'où 2.a).
\item La propriété 3) est claire puisque le complété de $A$ en $x$ est intègre donc
équidimensionnel.  \end{itemize} \end{proof} On verra plus bas (\ref{1patho-DVR}) qu'il existe de
nombreux anneaux de valuation discrète qui ne sont pas quasi-excellents.


\section{L'exemple de base : les anneaux locaux noethériens complets.} Expliquons avec  Nagata
pourquoi les anneaux locaux noethériens complets sont excellents\footnote{Ceci permet de construire
de nombreux exemples d'anneaux de valuation discrète excellents de caractéristique positive (par
complétion de schémas réguliers aux points de hauteur $1$).}.

\bigskip

La propriété 2.a) est tautologique. La formelle caténarité a été vue (\ref{1complet-formel-cat}).
Reste 1.a). Une  extension finie d'un anneau complet étant complet, on doit prouver le résultat
suivant (cf.~\egalong{IV}{2}{22.7.6}).

\begin{theoreme}[Nagata]\label{1complet-exc} Si $X$ est local noethérien intègre et complet
\footnote{D'après~(\ref{11->ouvert}), ceci entraîne que le lieu régulier d'un schéma local
noethérien complet est ouvert, qu'il soit intègre ou non.}, alors le lieu régulier est
ouvert.\end{theoreme} \begin{proof} On va distinguer les cas d'égales et d'inégales
caractéristiques.

\begin{itemize} \item[Cas I] (Cf.~\egalong{IV}{2}{0.22.7.6}.) Supposons que $A$ contienne un corps et
notons ${k_0}$ son corps premier (qui est parfait !) de sorte que le corps résiduel $k$ de $A_\p$
est séparable sur ${k_0}$ pour tout $\p\in\Spec(A)$. L'anneau $A_\p$ est régulier si et seulement si
$A_\p$ est formellement lisse sur ${k_0}$ (voir dans ce cas \egalong{IV}{2}{0.19.6.4}). D'autre part,
le théorème de structure de Cohen (\ref{1Co}) assure que $A$ est isomorphe à $k[[T_1,\cdots,T_n]]/I$
de sorte que $\p$ s'identifie à un idéal de $B=k[[T_1,\cdots,T_n]]$ contenant $I$. Le \emph{critère
jacobien de formelle lissité de Nagata} (\egalong{IV}{2}{0.22.7.3}) assure que $A_\p$ est régulier si
et seulement si il existe des ${k_0}$-dérivations $D_i,i=1,\cdots,m$ de $B$ dans $B$ et
$f_i,i=1,\cdots,m$ des éléments engendrant $I_\p$ tels que $\det(D_if_j)\not\in\p$. Cette condition
étant visiblement ouverte, le théorème suit.

\item[Cas II] Supposons que $A$ est d'inégale caractéristique, et donc de  corps des fractions $K$
de caractéristique nulle. D'après le théorème de structure de Cohen (\ref{1Co}), $A$ contient un
sous-anneau régulier (et complet) $B$ faisant de $A$ une $B$-algèbre de finie. Le corps des
fractions $L$ de $B$ est de caractéristique nulle comme $K$.  Quitte à remplacer $A$ par un localisé
$A[1/a]$, on peut supposer que $B$ est libre de rang fini sur $A$ de base $y_1,\cdots,y_m$.  Mais
$\Spec(A)\ra\Spec(B)$ est étale en dehors du fermé $d=\det_{A/B}(\mathrm{Tr}(y_iy_j))\not=0$ de
$\Spec(B)$, qui est non trivial puisqu'il contient le point générique, l'extension
$\Frac(B)/\Frac(A)$ étant séparable -- de caractéristique nulle -- ! Comme $B$ est régulier, le
théorème suit.  \end{itemize} \end{proof}


\begin{remarque}\label{1Coh-Gab} Ainsi, un anneau de séries formelles sur un corps est excellent.

Notons que la preuve se simplifie si on connaît l'amélioration de Gabber du théorème de structure de
Cohen (\refext{IV}{4-Cohen-Gabber égale} et \refext{IV}{4-Cohen-Gabber-Epp}) : si $A$ noethérien est
local complet et intègre, il contient un anneau $B$ isomorphe à un anneau de séries formelles sur un
anneau de Cohen ou un corps tel que $\Spec(A)\ra\Spec(B)$ est fini et \emph{génériquement étale}. On
n'a alors pas besoin de distinguer les caractéristiques des corps de fractions dans la preuve. Mais
la preuve de cette amélioration est difficile.  \end{remarque}


\section{Permanence par localisation et extension de type fini}\label{1inv-exc} La notion de (quasi)
excellence est remarquablement stable.  Précisément, on a


\begin{theoreme}\label{1th-permanent} Toute algèbre de type fini ou plus généralement
essentiellement de type fini sur un anneau excellent (resp.  quasi-excellent) est excellente (resp.
quasi-excellente). En particulier, tout localisé d'algèbre de type fini sur un corps ou sur un
anneau de Dedekind ($\ZZ$ par exemple) de corps des fractions de caractéristique nulle est
excellent.\end{theoreme}

(Rappelons que, dans ce contexte, un morphisme $\Spec(B) \ra \Spec(A)$
est dit \emph{essentiellement de type fini} si $B$ est une
localisation d'une $A$-algèbre de type fini par un système
multiplicatif.)

Expliquons les grandes lignes de la preuve.

\subsection{Condition 1)} Le passage au localisé ne pose pas de problème. Soit $B$ de type fini sur
$A$. Si $A$ vérifie 1.a) ou 1.b), les critères d'ouverture de Nagata (\egalong{IV}{2}{6.12.4 et
6.13.7}) entraînent qu'il en est de même de $B$.

\subsection{Condition 2)} C'est la partie la plus difficile de la théorie (\egalong{IV}{2}{7.4.4}),
entièrement due à Grothendieck. Le point le plus délicat est la localisation :

\begin{theoreme2} Si $A$ vérifie 2.a) (resp. 2.b) ou 2.c)), alors pour tout $\p\in\Spec(A)$,
l'anneau $A_\p$ vérifie 2.a) (resp. 2.b) ou 2.c)), autrement dit les fibres formelles en tout point
de $\Spec(A)$ sont géométriquement régulières (resp. géométriquement normales ou géométriquement
réduites).\end{theoreme2}

\begin{proof} La preuve se fait par réduction au cas complet. On se limite à la propriété 2.a), le
cas de 2.b) ou 2.c) se traitant de même. Soit $\m$ maximal contenant $\p\in\Spec(A)$ et soit $B$ le
complété $\m$-adique de $A$. Comme $A_\m\ra B$ est fidèlement plat, il existe $\q\in\Spec(B)$ au
dessus de $\p$. Par hypothèse, $A_\m\ra B$ est régulier. Les morphismes réguliers étant stables par
localisation, $A_\p\ra B_\q$ est régulier. On regarde alors le diagramme commutatif

$$\xymatrix{ \chap{A_\p}\ar[r]^{\hat f}&\chap{B_\q}\\ A_\p\ar[r]^f\ar[u]^\alpha&B_\q\ar[u]^\beta}.$$

Supposons que $\beta$ soit régulier. Alors, $\hat f\circ\alpha$ est régulier comme composé de deux
morphismes réguliers. Comme $\hat f$ est fidèlement plat (comme morphisme local  complété du
morphisme plat $f$), on déduit que $\alpha$ est régulier (exercice ou~\egalong{IV}{2}{6.6.1}) ce qu'on
voulait. On est donc ramené à $\beta$, donc au cas complet. La régularité de $\beta$ résulte alors
de

\begin{theoreme2}\label{1perm-fib-form-loc}Soit $B$ un anneau local noethérien complet $B$. Alors,
les fibres formelles de $B$ en $\q\in\Spec(B)$ sont géométriquement régulières.\end{theoreme2}

Ce théorème est le noyau dur de la théorie. On se ramène (\ref{1rem-fibre}) à étudier les fibres
formelles génériques. On montre donc dans un premier temps (\egalong{IV}{1}{0.22.3.3}) que si $\p$ est
un idéal premier de $A$  local noethérien complet intègre, la fibre formelle générique
$\chap{A_\p}\otimes_{A_\p}\Frac(A_\p)$  de $A_\p$ est formellement lisse sur $\Frac(A_\p)=\Frac(A)$
en tout point. Dans un second temps, on montre (\egalong{IV}{1}{0.22.5.8}) qu'une algèbre locale
noethérienne sur un corps est formellement lisse\footnote{Rappelons qu'une $k$-algèbre locale $B$
(muni de la topologie adique) est formellement lisse  sur $k$ si tout $k$-morphisme continu
d'algèbre $B\ra C/I$ avec $I^2=0$ se relève continûment à la $C$-algèbre discrète $C$.} si et
seulement si elle est géométriquement régulière\footnote{En fait, on n'a visiblement besoin que du
sens formellement lisse entraîne géométriquement régulier, qui est le plus facile. Notons que la
preuve de l'équivalence a été considérablement simplifiée par Faltings (\cite{Fal78} ou pour le
lecteur non germaniste \cite{CRT@Matsumura},~28.7).} \end{proof}


Une fois prouvée la permanence par localisation, on peut montrer :

\begin{theoreme2} Soit $B$ une $A$-algèbre de type fini. Si $A$ vérifie 2.a) (resp. 2.b) ou 2.c)),
alors $B$ vérifie 2.a) (resp. 2.b) ou 2.c)).\end{theoreme2}

La  preuve se fait par récurrence sur le nombre de générateurs de $B$. Grâce à l'invariance par
localisation, on se ramène aisément à l'étude des fibres formelles de $B$ en un idéal maximal dans
le cas où  $B$ engendré par \emph{un} élément et $A$ est complet. La preuve n'est pas facile, mais
beaucoup plus simple que celles de (\egalong{IV}{1}{0.22.3.3 et 0.22.5.8}).

\subsection{Condition 3).} De même que pour les conditions de type 2), la stabilité par localisation
et extension finie résulte comme plus haut (\egalong{IV}{2}{7.1.8}) du cas complet, la platitude du
morphisme de localisation permettant de descendre du complété à l'anneau -- ce n'est pas immédiat
malgré tout --. Le cas complet est facile comme on a vu (\ref{1reg-form-cat}).

\subsection{Application au cas local} Dans le cas local, la condition d'ouverture du lieu régulier
découle de 2.a). Précisons.

\begin{proposition2}\label{1loc-exc} \begin{itemize} \item [i)] Le lieu régulier d'un anneau local
noethérien vérifiant 2.a) est ouvert.  \item [ii)] En particulier, un anneau local noethérien est
quasi-excellent (resp.  excellent) si et seulement s'il vérifie 2.a) (resp.  s'il vérifie 2.a) et
3)).  \end{itemize} \end{proposition2}

\begin{proof}

Soit $f:X\ra Y$ un morphisme fidèlement plat de schémas noethériens à fibres régulières (resp.
normales ou réduites). Alors, $\OO_x$ est  régulier (resp.  normal ou réduit) si et seulement si
$\OO_{f(x)}$ l'est (\egalong{IV}{2}{6.4.2, 6.5.1}).  Notons $U_R(X)$ l'ensemble des $x\in X$ tel que
$R(\OO_x)$ est régulier (resp.  normal ou réduit). Autrement dit, on a $f^{-1}(U_R(X))=U_R(Y)$. Or
le lieu régulier ou normal d'un anneau complet intègre est ouvert  (\ref{1complet-exc}). De plus, le
morphisme de complétion d'un anneau local noethérien $A$ est régulier si et seulement si $A$ vérifie
2.a) d'après (\ref{1perm-fib-form-loc}). Il suit que 2.a) entraîne 1.a) (resp.  2.b) entraîne 1.b))
dans le cas local.  \end{proof}

\section{Comparaison avec ÉGA IV : le cas des anneaux universellement japonais}

Rappelons la définition usuelle des anneaux universellement japonais (\egalong{IV}{1}{0.23.1.1}).

\begin{definition}\label{1def-japonais} $X$  est dit \begin{enumerate} \item \emph{japonais} s'il
est intègre et si la clôture intégrale de $A$ dans toute extension finie\footnote{On peut se
contenter des extensions finies radicielles si l'on veut : exercice ou \egalong{IV}{1}{23.1.2}.} de son
corps des fractions est finie sur $A$\footnote{Comme module ou comme algèbre : c'est la même chose
car la clôture intégrale est entière sur $A$ par construction.}; \item \emph{universellement
japonais} si tout anneau intègre qui est extension de type fini de $A$ est japonais\footnote{Ou, ce
qui est équivalent (\egalong{IV}{2}{7.7.2}), si tout quotient intègre est japonais.}.  \end{enumerate}
\end{definition}


La définition d'anneau japonais n'est que technique en ce qu'elle ne sert qu'à définir la seule
notion véritablement utile (et vérifiable à vrai dire) : celle d'anneau universellement japonais.
Cette définition est compatible avec~\ref{1def-excellent}. Expliquons pourquoi. D'après Nagata, $X$
est universellement japonais au sens de~\ref{1def-japonais} si et seulement si $X$ vérifie 1.b) et
si tous les quotients intègres des localisés $\OO_{X,x}$ en les points fermés $x\in X$ sont japonais
(\egalong{IV}{2}{7.7.2}).  Or, le théorème de Zariski-Nagata (\egalong{IV}{2}{7.6.4}) assure  que les
quotients intègres de $\OO_{X,x}$ sont japonais si et seulement si les fibres formelles de
$\OO_{X,x}$  sont géométriquement réduites \footnote{\label{1Zar-Nag}Ou, de façon équivalente, que
le complété de toute $\OO_{X,x}$-algèbre finie et réduite est réduit. Comme d'habitude, la preuve se
fait par réduction au cas complet, et même régulier complet grâce au théorème de structure de Cohen.
Le caractère japonais de tels anneaux est garanti par le théorème de Nagata
(\egalong{IV}{1}{0.23.1.5}).}. D'où l'équivalence entre les deux définitions des anneaux
universellement japonais.



Si on renforce la condition 2.c) en 2.b) (fibres formelles géométriquement normales), le passage à
la clôture intégrale commute à la complétion. Précisément, on a (\egalong{IV}{2}{7.6.1 et 7.6.3})


\begin{proposition}\label{1prop-normal-versus-complet} Supposons que $A$ local noethérien vérifie
2.b) et soit réduit. Alors, la clôture intégrale $A'$ de $A$ dans son anneau total des fractions est
finie sur $A$ et son complété est isomorphe à la clôture intégrale de $\chap{A}$\footnote{Qui est
réduit puisque $A$ est japonais (cf.  note~(\ref{1Zar-Nag})).} dans son anneau total des fractions.
\end{proposition}

On déduit l'important critère d'intégrité du complété.

\begin{corollaire}\label{1-cor-normal-versus-complet} Soit $A$ local noethérien.  \begin{itemize}
\item[i)] Supposons $A$  intègre et vérifiant 2.b). Alors, $\chap{A}$ est intègre si et seulement si
$A$ est unibranche (\emph{i.e.} $A'$ local).  \item [ii)] Supposons $A$ hensélien. Alors $A$ est excellent
si et seulement s'il vérifie 2.a). Si $A$ est de plus intègre, il en est de même de son complété.
\end{itemize} \end{corollaire}

\begin{proof} Prouvons \textit{i)}. Comme $A$ est unibranche, la clôture intégrale $A'$ de $A$ est
locale : il en est de même de son complété $\chap{A'}$. D'après (\ref{1prop-normal-versus-complet}),
on a $\chap{A'}=(\chap{A})'$ et donc est normal. Or, un anneau normal et local est intègre. Comme
$\chap{A'}$ contient $\chap{A}$, le résultat suit.

Prouvons \textit{ii)}. D'après (\ref{1loc-exc}), on doit seulement se convaincre qu'un anneau local
hensélien vérifiant 2.a) vérifie aussi 3), \emph{i.e.} est formellement caténaire. On peut supposer $A$
intègre et on doit prouver que $\chap{A}$ est équidimensionnel. Mais comme $A$ est hensélien
intègre, il est unibranche~\egalong{IV}{4}{18.8.16}, donc $\chap{A}$ est intègre d'après le premier
point, ce qui assure l'équidimensionalité.\end{proof}

\section{Comparaison avec ÉGA IV : le cas des anneaux excellents}



La définition des anneaux noethériens excellents de Grothendieck est \textit{a priori} différente de
celle donnée ici. Notamment, elle fait intervenir, un peu bizarrement, l'\emph{universelle
caténarité} en lieu et place de la \emph{formelle caténarité}. Précisément, elle fait intervenir
trois propriétés. Dans cette partie $A$ désigne un anneau noethérien et $X=\Spec(A)$ le schéma
affine correspondant.

\begin{itemize} \item [1EGA)] Pour tout quotient intègre $B$ de $A$ et toute extension finie
radicielle $K'$ du corps des fractions $K$ de $B$, il existe une sous-$B$-algèbre finie $B'$ de $K'$
contenant $B$, de corps des fractions $K'$ telle que le lieu régulier de $\Spec(B')$ soit un ouvert
dense.  \item [2EGA)] Les fibres formelles de $X$ en tout point (fermé ou non) sont géométriquement
régulières.  \item [3EGA)] $A$ est universellement caténaire.  \end{itemize}

Les anneaux excellents au sens des \textbf{ÉGA} sont les anneaux noethériens vérifiant les 3
propriétés précédentes (\egalong{IV}{2}{7.8.2}).

Notons tout de suite, ce qui est élémentaire, que l'universelle caténarité de $A$ équivaut à celle
des anneaux locaux de $\OO_{X,x}$ en tous ses points fermés -- ou tous ses points si on préfère --
(\egalong{IV}{2}{5.6.3}).



Pour que la définition des anneaux excellents de Grothendieck (\egalong{IV}{2}{7.8.2}) soit la même que
(\ref{1def-excellent}), on doit prouver la proposition suivante.

\begin{proposition}\label{1equi-excellent}Pour tout anneau noethérien et $i=1,2,3$, les propriétés
$i$\emph{)} et $i$\emph{ÉGA)} son équivalentes. En particulier, les notions de quasi-excellence et
d'excellence de la première partie coïncident avec celles des \textbf{ÉGA}.\end{proposition}



\begin{proof}

La condition 1EGA) équivaut à 1) d'après~\egalong{IV}{2}{6.12.4} (seule la partie 1EGA) entraîne 1) est
délicate même si elle n'utilise pas le critère de régularité de Nagata mais seulement de l'algèbre
commutative standard -- essentiellement le critère de régularité par fibres et la non dégénérescence
de la trace des extensions finies séparables de corps --).

Pour l'équivalence de 2) et 2EGA), il faut se convaincre que la géométrique régularité des fibres
formelles en tout point fermé entraîne la géométrique régularité des fibres formelles en tout point
: c'est un cas particulier des propriétés de permanence (\ref{1inv-exc}).


Ceci prouve la compatibilité des définitions de la quasi-excellence.

Si $X$ vérifie 3), tous ses anneaux locaux sont formellement caténaires (permanence par
localisation, cf. la section \ref{1inv-exc}) et donc sont caténaires (\ref{1form-versus-cat}).
Comme tout schéma (affine) de type fini sur $X$ vérifie 3) (permanence par extension de type finie,
cf. la section \ref{1inv-exc}), on déduit que $X$ est universellement caténaire et donc $X$ vérifie
3EGA).


La réciproque est due à Ratliff :

\begin{proposition2}[Ratliff]\label{1Ratliff}\label{1form-cat->uni-cat} Un anneau noethérien
universellement caténaire est formellement caténaire.\end{proposition2}

Précisément, Ratliff prouve (\cite{Rat71}, 3.12) que si $A$ est caténaire, $A_\p$ est formellement
caténaire dès que $\p$ n'est pas maximal\footnote{Dans l'étrange terminologie de l'auteur, c'est la
condition $\text{depth}(\p)>0$, ce qui signifie donc que la \emph{dimension} de $A/\p$ est $>0$.}.
Pour montrer la proposition, on peut donc supposer $\p$ maximal et $A$ local intègre. Alors, $\p[X]$
est premier non maximal dans $A[X]$ de sorte que le complété $\p[X]$-adique $\chap{A[X]_{\p[X]}}$
est formellement équidimensionnel. Comme $\chap{A}\ra\chap{A[X]_{\p[X]}}$ est local et plat,
l'argument de platitude (\egalong{IV}{2}{7.1.3}) utilisé plus haut assure que $\chap{A}$ est
équidimensionnel.

\end{proof}


\section{Hensélisation et anneaux excellents}
\label{1-henselisation-anneaux-excellents}

Rappelons qu'un morphisme d'anneaux noethériens $A\ra B$ est dit absolument plat s'il est réduit à
fibres discrètes et si les extension résiduelles sont algébriques et séparables. Ou, de façon
équivalente, s'il est plat ainsi que le morphisme diagonal $B\otimes_AB\ra B$ (cf. \cite{Fer72},
prop. 4.1 et~\cite{Oli71}, 3.1). Lorsque $B$ est (localement) de type fini sur $A$, ceci équivaut au
fait que $B$ soit étale sur $A$. En particulier, les extensions résiduelles sont séparables de sorte
qu'un tel morphisme est en fait régulier. Par exemple, tout morphisme ind-étale est absolument plat.
On a alors le résultat suivant (\cite{Gre76}).


\begin{theoreme}\label{1th-greco} Soit $f:A\ra B$ un morphisme absolument plat d'anneaux
noethériens.  Alors \begin{itemize} \item[i)] Si $A$ est vérifie 2.a) (resp.  2.b) ou 2.c)), $B$ un
vérifie 2.a) (resp.  2.b) ou 2.c))\footnote{Ou plus généralement, si $A$ est un $\mathbf{P}$-anneau
au sens de Grothendieck (\egalong{IV}{2}{7.3}), $B$ est un $\mathbf{P}$-anneau}.  \item[ii)] Si $A$ est
universellement-japonais, $B$ est universellement japonais.  \item[iii)] Si $A$ est quasi-excellent,
$B$ est quasi-excellent.  \item[iv)] Si $A$ est excellent, $B$ est excellent.  \item[v)] Si $f$ est
fidèlement plat, la réciproque de i), ii) et iii) est vraie.  \item[vi)] Si $f$ est  fidèlement plat
et $B$ est localement intègre, la réciproque de iv) est vraie.  \end{itemize} \end{theoreme}


Comme les morphismes d'hensélisation et de stricte hensélisation sont absolument plats
(\egalong{IV}{4}{18.6.9 et 18.8.12}) et fidèlement plats (ils sont locaux), on trouve, en particulier,
que la quasi-excellence et l'excellence sont stables par hensélisation et hensélisation stricte et
que, le caractère quasi-excellent ou universellement japonais d'un anneau local se teste sur
l'hensélisé ou l'hensélisé strict. Dans le cas de l'hensélisé, ces résultats étaient connus de
Grothendieck (\egalong{IV}{4}{18.7}), notamment 18.7.6).

En revanche, on ne peut espérer une propriété de descente de l'excellence comme en vi) sans
condition d'intégrité locale (cf.\ref{1pathologie}).

\section{Complétion formelle et anneaux excellents}

Soit $I$ un idéal d'un anneau noethérien $A$ contenu dans son radical de Jacobson et $\gtilde A$ sa
complétion $I$-adique. On peut se demander si les propriétés d'excellence passent au complété. La
réponse est oui en général. Précisément, on a :

\begin{proposition} Soit $I$ un idéal d'un anneau noethérien $A$ contenu dans son radical de
Jacobson et $\gtilde A$ sa complétion $I$-adique.  \begin{itemize} \item [i)] Si $A$ est
(semi)-local quasi-excellent (resp. excellent), il en est de même de $\gtilde A$ ; \item [ii)] Si
$A$ est une $\QQ$-algèbre excellente, il en est de même de $\gtilde A$.  \end{itemize}
\end{proposition}


La permanence de la quasi-excellence dans le cas (semi)-local, \emph{i.e.} de la géométrique régularité des
fibres formelles, est due à Rotthaus (\cite{Rot80}), tandis que celle de l'universelle caténarité
est due à Seydi\footnote{Comme me l'a expliqué Christel Rotthaus (communication privée), si $A, B$
sont locaux tels que $A\subset B \subset \chap{A}$ et $\chap{B}=\chap{A}$, alors l'universelle
caténarité de $A$ entraîne celle de $B$.} (le théorème 1.12 de \cite{Sey70} prouve qu'un anneau de
série-formelles $A[[t_1,\cdots,t_n]]$ est universellement caténaire dès que $A$ l'est ; il suffit
alors de considérer des générateurs $i_1,\cdots,i_n$ de $I$ définissant une surjection
$A[[t_1,\cdots,t_n]]\twoheadrightarrow \gtilde A$). Pour \textit{ii)}, reste à étudier l'ouverture
du lieu régulier dans le.  C'est ce qui est fait dans \cite{Bro-Rot80}, en utilisant le théorème de
désingularisation d'Hironaka de façon cruciale. Les techniques de \cite{Bro-Rot80} ont d'ailleurs
permis de montrer que si le théorème de désingularisation était vrai dans le cas local excellent,
toute complétion $I$-adique d'anneau excellent comme plus haut serait excellente (\cite{Nis-Nis87}).

En fait, le résultat est général. Plus précisément, Gabber (\cite{Gab07}) peut remplacer le théorème
d'Hironaka par son théorème d'uniformisation (\refext{VII}{7-uniformisation}) dans les arguments de
\cite{Nis-Nis87} pour prouver le résultat suivant

\begin{theoreme}[Gabber]Soit $A$ un anneau noethérien $I$-adiquement complet.  Alors, si $A/I$ est
quasi-excellent, $A$ est quasi-excellent.\end{theoreme}

On ne peut remplacer quasi-excellent par excellent dans le théorème précédent. En effet, Greco
(\cite{Gre82}) a construit un idéal $I$ d'un anneau $A$ intègre de dimension $3$, noethérien
semi-local $I$-adiquement complet et séparé qui est quasi-excellent non excellent alors que $A/I$
est excellent. On peut même supposer que $A$ est une $\QQ$-algèbre. La construction se fait par
pincements d'idéaux maximaux de hauteurs différentes (cf.~\ref{1ex-non-uni-cat}). Malgré tout, comme
on vient de le voir, la formelle caténarité passe aux complétions partielles (\cite{Sey70}) de sorte
que le complété $I$-adique d'un anneau excellent $A$ est excellent dès lors que $I$ est contenu dans
le radical de Jacobson de $A$.

\section{Approximation d'Artin et anneaux excellents}

Rappelons la définition suivante (cf. \cite{Art69}).

\begin{definition}[M. Artin] Un anneau local noethérien $(A,\m)$ a la propriété d'approximation
(\emph{AP})  si pour toute variété affine $X$ de type fini sur $A$, l'ensemble $X(A)$ est dense dans
$X(\chap{A})$.  \end{definition}

Bien entendu, il revient au même de dire que pour tout $X$ comme plus haut, on a
$$X(\chap{A})\not=\emptyset\Rightarrow X(A)\not=\emptyset.$$

Si $A$ vérifie AP, $A$ est certainement hensélien. Mais l'exemple~\ref{1DVR-bad} prouve qu'il ne
suffit pas que $A$ soit hensélien pour qu'il possède la propriété d'approximation. En fait, Rotthaus
a observé que l'excellence était une condition nécessaire à l'approximation d'Artin :

\begin{lemme}[\cite{Rot90}] Un anneau local noethérien vérifiant AP est hensélien et excellent.
\end{lemme}

Soit $k$ un corps de caractéristique nulle muni d'une valuation non triviale à valeurs réelles.
Artin a prouvé (\cite{Art68}) que les anneaux de séries convergentes à coefficients dans $k$ (pas
nécessairement complets) ont la propriété d'approximation. Ils sont donc henséliens et
excellents.\footnote{Dans le même ordre d'idées, l'anneau $k\{x_1,\cdots,x_n\}$ des séries formelles
restreintes (séries dont la suite des coefficients tendent vers $0$) à coefficients dans un corps
valué complet non archimédien $k$ est excellent dès que $k$ est de caractéristique nulle ou que $k$
est de degré fini sur $k^p, p=\text{car}(k)$ (\cite{Gre-Val74}). Le cas général a été obtenu par
Kiehl (\cite{Kie69} et aussi \cite{Con99} pour une preuve et des développements). Ceci répond,
partiellement, à une question de Grothendieck (\egalong{IV}{2}{7.4.8 B)}). En revanche, si $k$ est
valué non archimédien non complet de caractéristique positive tel que le morphisme de complétion
$k\ra\chap{k}$ n'est pas séparable, Gabber sait prouver que $k\{x_1\}$ n'est pas excellent.}

La situation est maintenant complètement clarifiée grâce aux travaux de Popescu culminant avec le
résultat suivant (\cite{NPD@Swan}) :

\begin{theoreme}[Popescu] \label{1-Artin-Popescu}\begin{itemize} \item [i)]Soit $A\ra B$ un
morphisme   régulier d'anneaux noethériens. Alors, $B$ est limite inductive filtrante de
$A$-algèbres lisses.  \item [ii)] Tout anneau local noethérien hensélien et excellent satisfait la
propriété d'approximation \emph{AP}\footnotemark.  \end{itemize} \end{theoreme} \footnotetext{Voir
\cite{Spi99}, th. 11.3 pour un énoncé un peu plus général.}

Le fait que i) entraîne ii) est un simple exercice. En effet, si $A$ est quasi-excellent, le
morphisme de complétion $A\ra \chap{A}$ est régulier et donc on peut écrire $\chap{A}=\limind\ L$ où
$L$ est lisse sur $A$. Soit $X=\Spec(B)$ avec $B$ de type fini et $\hat a\in X(\chap{A})$ d'image
$\hat a(0)\in X(k)$ où $k$ corps résiduel de $A$.  Il existe donc $L$ lisse sur $A$ tel que $\hat a$
provienne de $l\in X(L)$. Comme $A$ est hensélien, il existe $a\in X(A)$ (tel que $a(0)=\hat
a(0))$\footnote{L'argument n'utilise que la géométrique régularité des fibres -- et le caractère
hensélien -- mais pas la formelle caténarité. Ce n'est pas paradoxal, car un anneau local hensélien
est excellent si et seulement si ses fibres formelles sont géométriquement régulières
(\ref{1loc-exc}).}.

\section{Exemples de méchants anneaux noethériens}\label{1pathologie} Soit $A$ un anneau noethérien
et $X=\Spec(A)$ le schéma affine correspondant.  Il ressortira de cet inventaire que les propriétés
désagréables des anneaux du point de vue de l'excellence n'ont en général pas seulement à voir avec
la caractéristique $>0$ mais peuvent aussi se produire pour des $\QQ$-algèbres.

\subsection{Formelle caténarité : condition 3)} \label{1ex-non-uni-cat} Regardons d'abord de mauvais
anneaux du point de vue de la formelle caténarité.  \subsubsection{La caténarité n'entraîne pas la
formelle caténarité.} Dans \egalong{IV}{2}{5.6.11}, Grothendieck construit un exemple d'anneau local
noethérien de dimension $2$ intègre, caténaire mais non universellement caténaire, donc non
formellement caténaire d'après~\ref{1form-cat->uni-cat} (\emph{i.e.} ne vérifiant pas 3)).


Expliquons la construction qui consiste à pincer une surface lisse sur un corps $k$, de
caractéristique nulle si l'on veut, le long de deux points de hauteurs différentes ayant des corps
résiduels isomorphes à $k$.


On part d'un corps $k$ extension transcendante pure de degré  infini sur son corps premier par
exemple, que l'on peut même supposer être $\QQ$. Soit $S$ une surface lisse munie d'un morphisme
projectif sur $T=\Spec(k[\tau])$ et $t\in S(T)$. On suppose qu'il existe un point $s\in S(k)$ de la
fibre $S_0$ de $S\ra T$ au dessus de $0\in T$ qui n'est pas dans l'image de $t$. Par exemple, on
peut prendre $S=\Spec(k[\sigma,\tau])$ avec $t$ la section d'image $\sigma=0$ et $s=(1,0)$. Les
corps $k(s)=k$ et $k(t)=\Frac(k[\tau])$ sont des extensions transcendantes pures de $\QQ$ de même
degré (infini) de sorte qu'on peut choisir un isomorphisme de corps $k(s)\isononcan k(t)$.  Ceci
permet de définir le sous-anneau $\OO_\Sigma$ de $\OO_S$ des fonctions qui coïncident en $s$ et $t$.
On dispose donc d'un morphisme $\pi:S\ra \Sigma$ qui envoie $s,t$ sur $\sigma\in\Sigma$.  Posons
$A=\OO_{\Sigma,\sigma}$ et soit $B$ l'anneau de coordonnées de
$S\times_\Sigma\Spec(\OO_{\Sigma,\sigma})$. Par construction, $\dim(B_s)=2$ et $\dim(B_t)=1$. Alors,
$A$ est noethérien, et $B$ est la normalisation de $A$ et est fini sur $A$. Comme $A$ est de
dimension $2$ et intègre il est évidemment caténaire. Si $A$ était universellement caténaire, la
formule de dimension (voir note~(\ref{1form-dim})) entraînerait $\dim(A)=\dim B_s=\dim B_t$, une
contradiction.

On peut même trouver pour tout $n\geq 2$ des anneaux locaux noethériens intègres de dimension $n$
vérifiant 2.a), caténaires non universellement caténaires donc ne vérifiant pas 3)
(\cite{Hei-Rot-Wie00}).

\subsubsection{La formelle caténarité ne se teste pas sur l'hensélisé.} Par des techniques de
pincements analogues de surfaces sur $k$, donc de caractéristique nulle si on veut, comme plus haut,
Grothendieck construit en effet un exemple d'anneau local non universellement caténaire (donc non
formellement caténaire) dont l'hensélisé  est excellent (\egalong{IV}{4}{18.7.7}). Quitte à changer de
base par la clôture séparable, on s'aperçoit que la formelle caténarité ne se teste pas plus sur
l'hensélisé strict, contrairement à la quasi-excellence (\ref{1th-greco}).

\subsubsection{La formelle caténarité n'entraîne certainement pas 2.a) (ni même 2.c)).} Par exemple,
un anneau de valuation discrète $A$ non excellent (cf. la partie~\ref{1patho-DVR}) a une fibre
formelle générique non géométriquement régulière (en effet, il est formellement caténaire
(\ref{1CM-form-cat}) et 2.a) entraîne 1.a) dans le cas local (\ref{1loc-exc})). Or, cette fibre
générique formelle est artinienne dans ce cas (elle ne contient pas l'idéal maximal de $\chap{A}$)
et donc la géométrique régularité équivaut ici à la géométrique réduction.


\subsubsection{Il existe des anneaux intègres \emph{normaux} non formellement caténaires.} Ogoma
(\cite{Ogo81}) a construit une $\QQ$-algèbre locale $A$ intègre \emph{normale} de dimension $3$ dont
le complété à une composante de dimension $2$ et une composante de dimension $3$ et donc n'est pas
équidimensionnel. Pire, cet anneau n'est même pas caténaire : il possède une infinité de chaînes
saturées d'idéaux premiers de longueur 2 ou 3.



\subsection{Quasi-excellence : conditions d'ouverture  1.a) et 1.b)}\label{1ex-non-1} On s'intéresse
ici à des anneaux ayant un lieu régulier ou normal non ouvert.

 Comme on le verra plus loin (\refext{XIX}{19-rem-cexgabber}), Gabber a construit un exemple de
schéma, qu'on peut même supposer être un $\QQ$-schéma,  intègre de dimension $1$, dont le lieu
régulier (ou normal, c'est la même chose ici) contient une infinité de points et en particulier
n'est pas ouvert. La construction assure que les fibres formelles sont géométriquement régulières.
Comme on est en dimension $1$, normalité et régularité coïncident de sorte qu'on a un exemple
vérifiant 2.a) et 3) mais pas 1.b).

 Dans~\cite{Rot79}, Rotthaus construit une $\QQ$-algèbre noethérienne locale intègre de dimension
$3$ qui est formellement caténaire, universellement japonaise mais dont le lieu régulier n'est pas
ouvert.

\subsection{Quasi-excellence. Fibres formelles : conditions 2a), 2b) et 2c)}\label{1ex-non-uni-2} On
s'intéresse ici à des anneaux ayant des fibres formelles non géométriquement régulières voire pire.

\begin{itemize} \item Rotthaus construit une $\QQ$-algèbre locale $A$ noethérienne de dimension $3$
régulière (donc formellement caténaire), universellement japonaise mais pas excellente
(\cite{Rot79}).  Précisément, la fibre formelle au dessus d'un point de hauteur $1$ n'est pas
régulière. Ainsi, elle vérifie 2.c), 3) car $A$ régulier mais pas 2.a).

Dans l'exemple d'Ogoma précédent, la fibre formelle générique est connexe non intègre (elle a une
composante de dimension $1$ et une de dimension $2$ qui se coupent), donc non normale. On a donc un
exemple de $\QQ$-algèbre locale (de dimension $3$) noethérienne (intègre et normale) ne vérifiant
pas 2.b).

On peut descendre d'une dimension : Nagata construit (\cite{Nag62}, ex. 7 de l'appendice A1) une
$\QQ$-algèbre locale $B$ qui est intègre normale, formellement caténaire et de dimension $2$ mais
dont le complété n'est pas intègre\footnote{La construction est la suivante : soient $x,y$
algébriquement indépendants sur $\QQ$ et $w=\sum_{i>0}a_ix^i\in\QQ[[x]]$ transcendant sur $K(x)$. On
pose $z_1=(y+w)^2$ et $z_{i+1}=(z-(y+\sum_{j<i}a_jx^j)^2)/x^i$. Soit $A$ le localisé de
$\QQ[x,y,z_i,i\geq 1]$ en $(x,y,z_i,i\geq 1)$. Alors, $B=A[X]/(X^2-z_1)$ est l'exemple cherché. On
vérifie facilement que la complétion de $A$ est $\QQ[[x,y]]$ de sorte que $\chap{B}=
\QQ[[x,y]][X]/(X^2-(y+w)^2)$ n'est pas intègre. Comme d'habitude dans ces constructions, c'est le
caractère noethérien de $A$ qui pose problème. Une fois ceci acquis, $A$ est régulier de dimension
$2$ et $B$ normal puisque que Cohen-Macaulay de dimension $2$ singulier uniquement à l'origine.
Notons que $B$ est formellement caténaire comme quotient d'un régulier.}.
D'après~\ref{1prop-normal-versus-complet}, ceci prouve que $B$ ne vérifie pas 2.b) (mais vérifie
3)).

\item En caractéristique $>0$, Rotthaus  construit également (\cite{Rot79}) une algèbre locale
noethérienne de dimension $2$ régulière universellement japonaise mais pas excellente. Dans ce
denier cas, comme les fibres formelles sont de dimension $<2$, elles ne sont pas non plus
géométriquement normales. Ainsi, elle vérifie 2.c), 3)  mais pas 2.b).

\item En caractéristique nulle, Ferrand et Raynaud ont construit (\cite{Fer-Ray70}, prop. 3.3 et 3.5)
une $\CC$-algèbre locale noethérienne $A$ intègre de dimension $2$ telle que

$\bullet$ le normalisé $A'$ de $A$ est l'anneau des séries convergentes noté $\CC\{x,y\}$ (ne pas
confondre avec l'hensélisé de $\CC[x,y]$) et donc est excellent.

$\bullet\ A$ n'est pas japonais (en fait, $A'$ n'est pas fini sur $A$).

$\bullet\ \chap{A}$ a des composantes immergées (de sorte que -- platitude -- la fibre formelle
générique a des composantes immergées et donc ne vérifie pas 2.c)).


$\bullet\ $ Le lieu normal de $A[[T]]$ n'est pas ouvert.  D'après~\egalong{IV}{2}{6.13.5} son lieu
régulier n'est donc pas ouvert non plus.

$\bullet$ L'anneau $A$ est formellement caténaire (le spectre de son complété est irréductible). Il
en est donc de même de $A[[T]]$ (\cite{Sey70}).


\item Dans~\cite{Nag62}, ex. 5 de l'appendice A1, Nagata construit même un anneau local noethérien
intègre de dimension $3$ (de caractéristique $>0$) dont la clôture intégrale n'est même pas
noethérienne ; en particulier, cet anneau n'est pas japonais\footnote{La construction est du même
type que celle d'un anneau de valuation discrète décrite dans la note (\ref{1note-DVR-bad}) dont on
reprend les notations. On considère cette fois-ci l'anneau $B=k^p[[X,Y,Z]][k]$ et
$d=Y\sum_{i>0}X_iX^i+Z\sum_{i>0}X_{2i+1}X^i$. L'anneau $B[d]$ convient.}.

\item Pire, à partir d'anneaux construits par Nagata, Seydi construit (\cite{Sey72}) un anneau
noethérien  intègre normal $A$ de dimension $3$ dont le corps des fractions est de caractéristique
nulle et dont le complété n'est pas réduit. En particulier, il est japonais mais pas universellement
japonais. Ogoma construit (\cite{Ogo81}) une $\QQ$-algèbre noethérienne normale, donc japonaise, qui
n'est ni universellement japonaise ni caténaire.

\item Les fibres formelles peuvent être épouvantables, même en dimension $1$ : Ferrand et Raynaud
construisent un $\CC$-schéma local intègre de dimension $1$ dont la fibre formelle générique est un
schéma artinien qui n'est même pas Gorenstein -- donc certainement non réduit -- (\cite{Fer-Ray70},
prop. 3.1) : $X$ ne vérifie pas 2.c). En particulier, $X$ n'est pas universellement japonais (et
donc pas quasi-excellent). Bien entendu, $X$ vérifie 1.a) et 3) pour des raisons de dimension.

\item  Les exemples d'anneaux de valuation discrète non excellents (donc de caractéristique
positive) donnent des exemples d'anneaux ne vérifiant pas 2.c) (2.a) et 2.c) sont équivalents en
dimension $\leq 1$) mais vérifiant 1.a) et 3).


\end{itemize}


\begin{remarque}\label{1DVR-bad}Nagata a construit (\cite{Nag62}, (E3.3)) un anneau de valuation
discrète dont la fibre formelle générique est une extension radicielle non triviale de son  corps
des fractions, donc non excellent\footnote{\label{1note-DVR-bad}Voici la construction : soit $k$ le
corps des fractions de $\FF_p[X_n,n> 0]$ et $K$ celui de $\chap{A}=k[[Y]]$. Soit $L$ le sous-corps
de $K=k((Y))$ corps des fractions de $A=k^p[[Y]][k]$. Le complété de $A$ est $\chap{A}$. On montre,
et c'est le point délicat, que $A$ est noethérien. L'outil est le critère de Cohen : un anneau
semi-local est noethérien si et seulement si les idéaux maximaux sont de type fini et les idéaux de
type fini fermés (\cite{Nag62}, 31.8). Son complété étant régulier, il est lui même régulier donc de
valuation discrète (dimension).  Soit $L$ le corps des fractions de $A$. On vérifie facilement que
$c=\sum_{n> 0}X_nY^n,$ n'est pas dans $L$. Choisissons une $p$-base $\{c_i\}$ de $K$ sur $L$
contenant $c$ (ce qui est possible car $c\not\in L^p$, cf.~\egalong{IV}{1}{21.4.3}). Soit $K_0$ le
corps engendré sur $L$ par les $c_i$ distincts de $c$. L'extension $K/K_0$ est radicielle de degré
$p$ par construction. L'anneau $A\cap K_0$ est un anneau de valuation discrète de complété $k[[Y]]$
de sorte que la fibre formelle générique n'est pas géométriquement réduite.}.  \end{remarque}

 On va voir maintenant que de tels anneaux se rencontrent très facilement.



\subsection{Méthode systématique de construction d'anneaux non quasi-excellents}\label{1patho-DVR}
En fait, on peut construire (Orgogozo) de façon systématique de très nombreux anneaux de valuation
discrète non quasi-excellents. Précisons\footnote{Cette construction généralise en fait, de façon
indépendante, un exemple obtenu par Rotthaus dans~\cite{Rot97}}.


\begin{proposition}\label{1lot-bad-DVR}Soit $k((t))$ le corps des séries de Laurent à coefficients
dans un corps $k$ de caractéristique $p>0$ muni de sa valuation $t$-adique et $L/k$ une
sous-extension de type fini de $k((t))/k$ de degré de transcendance $>1$ sur $k$. Alors, le
sous-anneau $A$ de $L$ des éléments de valuation $\geq 0$ est un anneau de valuation discrète non
excellent.  \end{proposition}

\begin{proof}

Soit ${\bar L}$ un corps de caractéristique $>0$. Le $p$-rang est la dimension, finie ou non, de
$\Omega_{\bar L}$, le module des différentielles absolues. C'est aussi $\log_p([{\bar L}:{\bar
L}^p])$ où $[{\bar L}:{\bar L}^p]$ est la dimension de ${\bar L}$ sur ${\bar L}^p$
(\egalong{IV}{1}{21.3.5}).  La remarque clef est que le $p$-rang croit par extension de corps
\emph{séparable} ${\bar K}/{\bar L}$,  finie ou non, puisqu'on a dans ce cas un plongement ${\bar
K}\otimes_{\bar L}\Omega_{\bar L}\hookrightarrow \Omega_{{\bar K}}$ (\egalong{IV}{2}{20.6.3})
\begin{equation}\label{1ineq-sep} [{\bar L}:{\bar L}^p]\leq [{\bar K}:{\bar K}^p].  \end{equation}

Par ailleurs, si ${\bar L}$ est de \emph{type fini} sur un corps $k$, on a (\cite{Bou50}, V.6.3)

\begin{equation}\label{1prang} [{\bar L}:{\bar L}^p]=p^{\text{degtr}_k({\bar L})}[k:k^p].
\end{equation}


Plaçons nous dans la situation du lemme. L'anneau $A$ est de valuation discrète par construction et
son corps des fractions est $L$. L'hensélisé $A^h$ est local régulier de dimension $1$ donc intègre
et son corps des fractions $K=\Frac(A^h)$ contient $L=\Frac(A)$.  Le  complété $\chap{A^h}$ est un
anneau de séries formelles $\chap{K}=k[[\varpi]],\varpi $ uniformisante de $A^h$ (comme complété
d'une $k$-algèbre locale régulière de dimension $1$) et son corps des fractions  $\chap{K}$ est la
fibre générique formelle de $\Spec(\chap{A^h})\ra\Spec(A^h)$.

Supposons que $A^h$ soit quasi-excellent (précisément vérifie 2.a) de sorte que l'extension
$\chap{K}/K$ est \emph{séparable}.


On a donc dans ce cas $$[K:K^p]\leq[\chap{K}:\chap{K}^p].$$ Comme $\chap{K}=k((\varpi)),$ on a
$$[\chap{K}:\chap{K}^p]=p[k:k^p].$$

On a donc (\ref{1ineq-sep}), l'extension $K/L$ étant séparable, $$[K:K^p]\geq [L:L^p].$$ de sorte
que, grâce à (\ref{1prang}), on a $$p[k:k^p]=[\chap{K}:\chap{K}^p]\geq [L:L^p]>p[k:k^p],$$ une
contradiction. Ceci interdit à $A$ également d'être quasi-excellent (\ref{1th-greco}).\end{proof}

\ifx\danslelivre\undefined
\bibliography{../fichiers-configuration-latex/bibliographie}
\bibliographystyle{../fichiers-configuration-latex/smfalpha}
\end{document}
\fi
